%-----------------------------------------------------------------------------
% Załącznik II
%-----------------------------------------------------------------------------

\cleardoublepage
\section*{II. Poziemscy spod Łochowa}
\label{sec:poziemscy}

Prapoczątków rodziny Poziemskich należałoby szukać w~okolicach miasta Łochów
położonego na~środkowym Mazowszu, mniej więcej 30~kilometrów
na~północny-zachód od~Węgrowa oraz 70~kilometrów na północny-wschód
od~Warszawy. Konkretnie pochodzili oni prawdopodobnie ze~wsi Budy Łosiewickie
(obecnie część wsi Łosiewice), gdzie w~1806~roku urodził się Andrzej
Poziemski\indeksosoby{Andrzej}{Poziemski}{1806}{Budy Łosiewickie}
-- pierwszy członek rodziny Poziemskich, który około 1825~roku przeprowadził
się w~okolice Okuniewa, a~dokładnie do~wsi Kąty-Majdan (obecnie część wsi
Majdan w~powiecie wołomińskim).

\begin{figure}[!ht]
    \centering
    % Pierwszy wypis
    \includegraphics[width=1.0\linewidth]{
        1801_Franciszek_Poziemski_Anna_Burakowska_akt_ślubu_parafia_Kamonionna_37.png}
    
    \vspace{0.1cm} % Odstęp między wypisami

    % Drugi wypis
    \includegraphics[width=0.82\linewidth]{
        1806_Andrzej_Poziemski_akt_chrztu_parafia_Kamonionna_266.png}
    
    % Wspólny podpis dla obu wypisów
    \captionsetup{format=hang}
    \caption{Fragmenty wypisów z~sumariusza alfabetycznego dla ksiąg
    urodzonych, zaślubionych, zmarłych z~lat 1766--1810 w~parafii Kamionna --
    na~czerwono zaznaczono wpisy odnoszące się do rodziny Poziemskich
    \cite{par_kamionna0}.}
    \label{fig:fpoziemski_1801}
\end{figure}

\begin{figure}[!ht]
    \centering \includegraphics[width=1.0\linewidth]{
        1828_Andrzej_Poziemski_Rozalia_Rosik_akt_ślubu_parafia_Okuniew_6.JPG}
    \captionsetup{format=hang}
    \caption{Akt ślubu Andrzeja Poziemskiego i~Rozali Rosik (z~domu Rudnik) --
    par. Okuniew 1828~rok (6/1828) \cite{par_okuniew0}.}
    \label{fig:apoziemski_1828}
\end{figure}

Najstarsze wzmianki o~rodzinie Poziemskich pochodzą z~początku XIX~wieku z~akt
metrykalnych parafii rzymskokatolickiej pod wezwaniem Niepokalanego Poczęcia
Najświętszej Maryi Panny w~Kamionnej, do~której należały wówczas Budy
Łosiewickie. Według sumariusza alfabetycznego dla ksiąg urodzonych,
zaślubionych, zmarłych z~lat 1766--1810 \cite{par_kamionna0} w~tej parafii,
w~1801~roku Franciszek
Poziemski\indeksosoby{Franciszek}{Poziemski}{?}{Budy Łosiewickie} wziął ślub
z~Anną Borakowską\indeksosoby[Borakowska]{Anna}{Poziemska}{?}{Budy Łosiewickie},
a~w~1806~roku urodził im~się syn
Andrzej\indeksosoby{Andrzej}{Poziemski}{1806}{Budy Łosiewickie}
(patrz ryc. \ref{fig:fpoziemski_1801}). Nie udało mi~się niestety dotrzeć
do~dokładnej treści tych aktów metrykalnych, więc nie dysponuję innymi
szczegółami odnośnie tych zdarzeń. Jedynym niezależnym źródłem, które
potwierdza te~informacje jest akt ślubu Andrzeja -- dnia 17~lutego 1828~roku
ożenił się on~w~parafii świętego Stanisława Kostki w~Okuniewie z~Rozalią
Rosik\indeksosoby[Rudnik][Rosik]{Rozalia}{Poziemska}{1800}{Ossów}
(ur.~1800~r.~-- zm.~?) wdową po~Filipie Rosiku, córką Franciszka i~Franciszki
Rudników, urodzoną we~wsi Ossów pod Wołominem. Z~aktu tego wynika,
iż~w~momencie ślubu Andrzej\indeksosoby{Andrzej}{Poziemski}{1806}{Budy Łosiewickie}
miał dwadzieścia jeden lat, co~potwierdzałoby jego narodziny w~okolicach
1806~roku, jego rodzicami byli
Franciszek\indeksosoby{Franciszek}{Poziemski}{?}{Budy Łosiewickie} oraz
Anna\indeksosoby[Borakowska]{Anna}{Poziemska}{?}{Budy Łosiewickie} Poziemscy
oraz że~urodził się on~w~miejscowości Budy w~parafii Kamionna (patrz ryc.
\ref{fig:apoziemski_1828}) -- wszystko wskazuje na~to, iż~chodzi właśnie
o~Budy Łosiewickie. Dokładny zapis tego aktu jest następujący (pisownia
oryginalna):

\begin{centeredquote}
\textit{... między Andrzeiem Poziemskim młodzianem, mającym lat dwadzieścia
ieden, urodzonym w~Budach Parafii Kamionna, Synem niegdyś Franciszka iuż
zmarłego i~Anny żyjącej Poziemskich Małżonków, utrzymującym się ze~służby
i~zamieszkałym na Kącie Maydan w~Parafii Cygowskiey ...}
\end{centeredquote}

Z~zapisów aktu ślubu małżonków Poziemskich wynika również,
iż~Andrzej\indeksosoby{Andrzej}{Poziemski}{1806}{Budy Łosiewickie} przed
wstąpieniem w~związek małżeński utrzymywał się ze~służby i~mieszkał we~wsi
Kąty-Majdan należącej do~parafii Cygów\footnote{Obecnie parafia świętego Jana
Chrzciciela i~świetego Wojciecha Biskupa Męczennika w~Poświętnem.}.

\begin{figure}[!ht]
    \centering \includegraphics[width=1.0\linewidth]{
        1851_Franciszek_Poziemski_Magdalena_Grzybowska_akt_ślubu_parafia_Okuniew_wpis_24.jpg}
    \captionsetup{format=hang}
    \caption{Akt ślubu Franciszka Poziemskiego i~Magdaleny Grzybowskiej --
    par. Okuniew 1851~rok (24/1851) \cite{par_okuniew1}.}
    \label{fig:fpoziemski_1851}
\end{figure}

Ze~względu na~brak dostępu do~kompletu akt metrykalnych z~parafii Okuniew
sprzed 1846~roku, nie wiemy ilu dzieci dochowali się małżonkowie
Andrzej\indeksosoby{Andrzej}{Poziemski}{1806}{Budy Łosiewickie}
i~Rozalia\indeksosoby[Rudnik][Rosik]{Rozalia}{Poziemska}{1800}{Ossów}
Poziemscy, wiemy natomiast, że~około 1829~roku w~Budziskach urodził
im~się syn
Franciszek\indeksosoby{Franciszek}{Poziemski}{1829}{Budziska} (ur.~1829~r.~--
zm.~1881~r.). Z~aktu jego ślubu z~Magdaleną
Grzybowską\indeksosoby{Magdalena}{Grzybowska}{1827}{Kaleń} z~1851~roku wynika,
iż Andrzej i~Rozalia Poziemscy utrzymywali się z~pracy jako służący
w~charakterze pasterzy w~Długiej Szlacheckiej (patrz ryc.
\ref{fig:fpoziemski_1851}) -- dokładny zapis aktu określa ich jako (pisownia
oryginalna):

\begin{centeredquote}
\textit{... Pasterzy w~Długiej Szlacheckiey zamieszkałych ...}
\end{centeredquote}

Dnia 23~listopada 1851~roku dwudziestojednoletni Franciszek
Poziemski\indeksosoby{Franciszek}{Poziemski}{1829}{Budziska} ożenił
się z~dwudziestopięcioletnią Magdaleną
Grzybowską\indeksosoby{Magdalena}{Grzybowska}{1827}{Kaleń} (ur.~1827~r.~--
zm.~1897~r.), córką nieżyjących już Wojciecha i~Agaty (z~domu Gumienny)
Grzybowskich, urodzoną we~wsi Kaleń\footnote{Wieś ta~obecnie nie istnieje --
została zlikwidowana wraz z~utworzeniem poligonu wojskowego w~Rembertowie pod
koniec XIX~wieku.} należącej do~parafii Okuniew (patrz ryc.
\ref{fig:fpoziemski_1851}). Zarówno Franciszek, jak i~Magdalena, przed
ślubem utrzymywali się ze~służby i~oboje zamieszkiwali we~wsi Budziska
(pisownia oryginalna):

\begin{centeredquote}
\textit{... małżeństwo między Franciszkiem Poziemskim ... utrzymującym się
ze~służby zamieszkałym w~Budziskach, a~Panną Magdaleną Grzybowską ...
utrzymującą się ze~służby, w~Budziskach zamieszkałą ...}
\end{centeredquote}

\begin{figure}[!ht]
    \centering \includegraphics[width=1.0\linewidth]{
        1856_Jan_Poziemski_akt_chrztu_parafia_Okuniew_wpis_32.jpg}
    \captionsetup{format=hang}
    \caption{Akt chrztu Jana Poziemskiego -- par. Okuniew 1856~rok (32/1856)
    \cite{par_okuniew1}.}
    \label{fig:jpoziemski_1856}
\end{figure}

Para dochowała się co~najmniej trzech synów:
Jana\indeksosoby{Jan}{Poziemski}{1856}{Budziska} (ur.~1856~r.~-- zm.~1862~r.),
Antoniego\indeksosoby{Antoni}{Poziemski}{1864}{Cisie} (ur.~1864~r.~--
zm.~1910~r.) oraz Piotra\indeksosoby{Piotr}{Poziemski}{1867}{Cisie}
(ur.~1867~r.~-- zm.~1921~r.). Z~aktu chrztu ich najstarszego syna Jana wynika,
iż~po~ślubie
Franciszek\indeksosoby{Franciszek}{Poziemski}{1829}{Budziska}
i~Magdalena\indeksosoby{Magdalena}{Grzybowska}{1827}{Kaleń} Poziemscy
zamieszkali prawdopodobnie w~pobliskiej do~Budzisk wsi Zabraniec, gdzie
Franciszek pracował w~charakterze wyrobnika (patrz ryc.
\ref{fig:jpoziemski_1856}).

\begin{figure}[!ht]
    \centering \includegraphics[width=1.0\linewidth]{
        1862_Jan_Poziemski_akt_zgonu_parafia_Długa_Kościelna_wpis_28.jpg}
    \captionsetup{format=hang}
    \caption{Akt zgonu Jana Poziemskiego -- par. Długa Kościelna 1862~rok
    (28/1862) \cite{par_dluga}.}
    \label{fig:jpoziemski_1862}
\end{figure}

Około 1860~roku rodzina Poziemskich przeprowadziła się do~oddalonej o~około
10~kilometrów na południe od~Zabrańca wsi Cisie -- wskazują na~to~zapisy aktu
zgonu ich sześcioletniego syna
Jana\indeksosoby{Jan}{Poziemski}{1856}{Budziska}. W~Cisiu również urodzili się
pozostali synowie małżeństwa Poziemskich, 14~stycznia 1864~roku na~świat
przyszedł Antoni\indeksosoby{Antoni}{Poziemski}{1864}{Cisie} (patrz ryc.
\ref{fig:apoziemski_1864}), a~31~lipca 1867~roku
Piotr\indeksosoby{Piotr}{Poziemski}{1867}{Cisie} (patrz ryc.
\ref{fig:ppoziemski_1867}). Akta chrztu Antoniego i~Piotra wskazują, iż~ich
ojciec Franciszek\indeksosoby{Franciszek}{Poziemski}{1829}{Budziska}, podobnie
jak w~Zabrańcu, trudnił się w~Cisiu pracą w~charakterze wyrobnika / parobka.

\begin{figure}[!ht]
    \centering \includegraphics[width=1.0\linewidth]{
        1864_Antoni_Poziemski_akt_chrztu_parafia_Długa_Kościelna_wpis_8.jpg}
    \captionsetup{format=hang}
    \caption{Akt chrztu Antoniego Poziemskiego -- par. Długa Kościelna
    1864~rok (8/1864) \cite{par_dluga}.}
    \label{fig:apoziemski_1864}
\end{figure}

\begin{figure}[!ht]
    \centering \includegraphics[width=0.9\linewidth]{
        1867_Piotr_Poziemski_akt_chrztu_parafia_Długa_Kościelna_wpis_71.jpg}
    \captionsetup{format=hang}
    \caption{Akt chrztu Piotra Poziemskiego -- par. Długa Kościelna 1867~rok
    (71/1867) \cite{par_dluga}.}
    \label{fig:ppoziemski_1867}
\end{figure}

Prawdopodobnie w~okolicach 1870~roku
Franciszek\indeksosoby{Franciszek}{Poziemski}{1829}{Budziska}
i~Magdalena\indeksosoby{Magdalena}{Grzybowska}{1827}{Kaleń} Poziemscy wraz
z~synami Antonim\indeksosoby{Antoni}{Poziemski}{1864}{Cisie}
i~Piotrem\indeksosoby{Piotr}{Poziemski}{1867}{Cisie} przeprowadzili się
do~pobliskiej do~Cisia wsi Chobot, gdzie osiedlili się już na~stałe. Wskazują
na~to~zapisy aktów zgonu Franciszka z~10~kwietnia 1881~roku (patrz ryc.
\ref{fig:fpoziemski_1881}) oraz Magdaleny z~14~marca 1897~roku (patrz ryc.
\ref{fig:mpoziemska_1897}) -- oboje zmarli właśnie we~wsi Chobot.

\begin{figure}[!ht]
    \centering \includegraphics[width=0.9\linewidth]{
        1881_Franciszek_Poziemski_akt_zgonu_parafia_Długa_Kościelna_wpis_24.jpg}
    \captionsetup{format=hang}
    \caption{Akt zgonu Franciszka Poziemskiego -- par. Długa Kościelna
    1881~rok (24/1881) \cite{par_dluga}.}
    \label{fig:fpoziemski_1881}
\end{figure}

\begin{figure}[!ht]
    \centering \includegraphics[width=1.0\linewidth]{
        1897_Magdalena_Poziemska_Grzybowska_akt_zgonu_parafia_Długa_Kościelna_wpis_22.jpg}
    \captionsetup{format=hang}
    \caption{Akt zgonu Magdaleny Poziemskiej z~domu Grzybowskiej -- par.
    Długa Kościelna 1897~rok (22/1897) \cite{par_dluga}.}
    \label{fig:mpoziemska_1897}
\end{figure}

\begin{figure}[!ht]
    \centering \includegraphics[width=1.0\linewidth]{
        1864_Józefa_Tomaszewska_akt_chrztu_parafia_Długa_Kościelna_wpis_98.jpg}
    \captionsetup{format=hang}
    \caption{Akt chrztu Józefy Tomaszewskiej -- par. Długa Kościelna 1864~rok
    (98/1864) \cite{par_dluga}.}
    \label{fig:jtomaszewska_1864}
\end{figure}

\begin{figure}[!ht]
    \centering \includegraphics[width=0.93\linewidth]{
        1886_Antoni_Poziemski_Józefa_Tomaszewska_akt_ślubu_parafia_Długa_Kościelna_wpis_21.jpg}
    \captionsetup{format=hang}
    \caption{Akt ślubu Antoniego Poziemskiego i~Józefy Tomaszewskiej --
    par. Długa Kościelna 1886~rok (21/1886) \cite{par_dluga}.}
    \label{fig:apoziemski_1886}
\end{figure}

Z~perspektywy historii rodziny Gomulskich z~Desna najistotniejsze wydają się
dalsze losy starszego z~synów
Franciszka\indeksosoby{Franciszek}{Poziemski}{1829}{Budziska}
i~Magdaleny\indeksosoby{Magdalena}{Grzybowska}{1827}{Kaleń} Poziemskich -
Antoniego\indeksosoby{Antoni}{Poziemski}{1864}{Cisie}\footnote{Ich młodszy
syn Piotr\indeksosoby{Piotr}{Poziemski}{1867}{Cisie} ożenił się w~1889~roku
z~Małgorzatą Zagańczyk\indeksosoby[Zagańczyk]{Małgorzata}{Poziemska}{1865}{Rysie}
i~dalsze losy swojej rodziny związał z~wsią~Chobot.}. Dnia 18~maja
1886~roku w~parafii świętej Anny w~Długiej Kościelnej Antoni Poziemski wziął
ślub ze~swoją rówieśniczką Józefą
Tomaszewską\indeksosoby[Tomaszewska]{Józefa}{Poziemska}{1864}{Chobot}
(ur.~1864~r.~-- zm.~?) panną urodzoną i~mieszkającą, tak jak Antoni, we~wsi
Chobot, córką Szymona i~Józefy (z~domu Przybyszewskiej) Tomaszewskich
(patrz ryc. \ref{fig:apoziemski_1886}). Małżonkowie Poziemscy dochowali się
wspólnie co~najmniej siódemki dzieci:

\begin{itemize}
    \item Franciszki\indeksosoby{Franciszka}{Poziemska}{1889}{Chobot}
    (ur.~1889~r.~-- zm.~1892~r.) - patrz ryc. \ref{fig:fpoziemska_1889} oraz
    \ref{fig:fpoziemska_1892},
    \item Aleksandra\indeksosoby{Aleksander}{Poziemski}{1892}{Chobot}
    (ur.~1892~r.~-- zm.~1967~r.) - patrz ryc. \ref{fig:apoziemski_1892},
    \item Franciszka\indeksosoby{Franciszek}{Poziemski}{1894}{Chobot}
    (ur.~1894~r.~-- zm.~1917~r.)  - patrz ryc. \ref{fig:fpoziemski_1894},
    \item Stefanii\indeksosoby[Poziemska]{Stefania}{Osica}{1897}{Chobot}
    (ur.~1897~r.~-- zm.~?) - patrz ryc. \ref{fig:spoziemska_1897},
    \item Michała\indeksosoby{Michał}{Poziemski}{1901}{Chobot}\footnote{Michał
    Poziemski był pradziadkiem autora niniejszej książki.} (ur.~1901~r.~--
    zm.~1947~r.) - patrz ryc. \ref{fig:mpoziemski_1901},
    \item Bezimiennego\indeksosoby{Bezimienny}{Poziemski}{1904}{Desno}
    (ur.~1904~r.~-- zm.~1904~r.) - patrz ryc. \ref{fig:bpoziemski_1904},
    \item Czesławy\indeksosoby{Czesława}{Poziemska}{1906}{Desno}
    (ur.~1906~r.~-- zm.~1922~r.) - patrz ryc. \ref{fig:czpoziemska_1906}.
\end{itemize}

\begin{figure}[!ht]
    \centering \includegraphics[width=0.85\linewidth]{
        1889_Franciszka_Poziemska_akt_chrztu_parafia_Długa_Kościelna_wpis_106.jpg}
    \captionsetup{format=hang}
    \caption{Akt chrztu Franciszki Poziemskiej -- par. Długa Kościelna
    1889~rok (106/1889) \cite{par_dluga}.}
    \label{fig:fpoziemska_1889}
\end{figure}

\begin{figure}[!ht]
    \centering \includegraphics[width=0.87\linewidth]{
        1892_Franciszka_Poziemska_akt_zgonu_parafia_Długa_Kościelna_wpis_18.jpg}
    \captionsetup{format=hang}
    \caption{Akt zgonu Franciszki Poziemskiej -- par. Długa Kościelna
    1892~rok (18/1892) \cite{par_dluga}.}
    \label{fig:fpoziemska_1892}
\end{figure}

\begin{figure}[!ht]
    \centering \includegraphics[width=0.87\linewidth]{
        1892_Aleksander_Poziemski_akt_chrztu_parafia_Długa_Kościelna_wpis_15.jpg}
    \captionsetup{format=hang}
    \caption{Akt chrztu Aleksandra Poziemskiego -- par. Długa Kościelna
    1892~rok (15/1892) \cite{par_dluga}.}
    \label{fig:apoziemski_1892}
\end{figure}

\begin{figure}[!ht]
    \centering \includegraphics[width=0.81\linewidth]{
        1894_Franciszek_Poziemski_akt_chrztu_parafia_Długa_Kościelna_wpis_123.jpg}
    \captionsetup{format=hang}
    \caption{Akt chrztu Franciszka Poziemskiego -- par. Długa Kościelna
    1894~rok (123/1894) \cite{par_dluga}.}
    \label{fig:fpoziemski_1894}
\end{figure}

\begin{figure}[!ht]
    \centering \includegraphics[width=0.81\linewidth]{
        1897_Stefania_Poziemska_akt_chrztu_parafia_Długa_Kościelna_wpis_48.jpg}
    \captionsetup{format=hang}
    \caption{Akt chrztu Stefanii Poziemskiej -- par. Długa Kościelna
    1897~rok (48/1897) \cite{par_dluga}.}
    \label{fig:spoziemska_1897}
\end{figure}

\begin{figure}[!ht]
    \centering \includegraphics[width=0.89\linewidth]{
        1901_Michał_Poziemski_akt_chrztu_parafia_Długa_Kościelna_wpis_106.jpg}
    \captionsetup{format=hang}
    \caption{Akt chrztu Michała Poziemskiego -- par. Długa Kościelna
    1901~rok (106/1901) \cite{par_dluga}.}
    \label{fig:mpoziemski_1901}
\end{figure}

\begin{figure}[!ht]
    \centering \includegraphics[width=0.89\linewidth]{
        1904_Bezimienny_Poziemski_akt_zgonu_parafia_Długa_Kościelna_wpis_10.jpg}
    \captionsetup{format=hang}
    \caption{Akt zgonu Bezimiennego Poziemskiego -- par. Długa Kościelna
    1904~rok (10/1904) \cite{par_dluga}.}
    \label{fig:bpoziemski_1904}
\end{figure}

\begin{figure}[!ht]
    \centering \includegraphics[width=0.97\linewidth]{
        1906_Czesława_Poziemska_akt_chrztu_parafia_Długa_Kościelna_wpis_123.jpg}
    \captionsetup{format=hang}
    \caption{Akt chrztu Czesławy Poziemskiej -- par. Długa Kościelna
    1906~rok (123/1906) \cite{par_dluga}.}
    \label{fig:czpoziemska_1906}
\end{figure}

\begin{figure}[!ht]
    \centering \includegraphics[width=0.97\linewidth]{
        1910_Antoni_Poziemski_akt_zgonu_parafia_Długa_Kościelna_wpis_66.jpg}
    \captionsetup{format=hang}
    \caption{Akt zgonu Antoniego Poziemskiego -- par. Długa Kościelna
    1910~rok (66/1911) \cite{par_dluga}.}
    \label{fig:apoziemski_1910}
\end{figure}

Antoni\indeksosoby{Antoni}{Poziemski}{1864}{Cisie}
i~Józefa\indeksosoby[Tomaszewska]{Józefa}{Poziemska}{1864}{Chobot} Poziemscy
zamieszkali po~ślubie w~swojej rodzinnej wsi Chobot, gdzie urodziła się piątka
ich najstarszych dzieci, po~czym około~1905~roku przeprowadzili się
do~sąsiedniej wsi \emph{Desna}\footnote{Antoni Poziemski był jednym
ze~współudziałowców spółki \emph{Desna} -- 13~marca 1905~roku odkupił
od~Bronisława Piotrowskiego\indeksosoby{Bronisław}{Piotrowski}{1868}{Wola
Rafałowska}  udziały w~tej spółce, co~opisałem szerzej w~rozdziale
V~niniejszej książki.}, gdzie przyszła na~świat dwójka ich pozostałych dzieci.
Antoni\indeksosoby{Antoni}{Poziemski}{1864}{Cisie} zmarł we~wsi \emph{Desna}
dnia 1~października 1910~roku, w~wieku 46~lat (patrz ryc.
\ref{fig:apoziemski_1910}). Pozostawił po~sobie owdowiałą żonę 
Józefę\indeksosoby[Tomaszewska]{Józefa}{Poziemska}{1864}{Chobot}\footnote{Niestety
data zgonu Józefy Poziemskiej nie jest znana autorowi niniejszej książki,
na~pewno dożyła ona do~1919~roku, gdyż została wymieniona jako osoba żyjąca
w~akcie ślubu swojego najstarszego syna
Aleksandra\indeksosoby{Aleksander}{Poziemski}{1892}{Chobot}.} oraz
piątkę\footnote{Najstarsza córka Antoniego i~Józefy Poziemskich,
Franciszka\indeksosoby{Franciszka}{Poziemska}{1889}{Chobot}, zmarła w~1892~roku
w~wieku niespełna trzech lat, a~ich najmłodszy syn przyszedł na~świat martwy,
toteż nie zostało mu~nadane nawet imię -- takie osoby określam konsekwentnie
w~tej książce jako \emph{Bezimienny} / \emph{Bezimienna}.} nieletnich dzieci.
Spośród tej piątki tylko troje założyło swoje własne rodziny:
Aleksander\indeksosoby{Aleksander}{Poziemski}{1892}{Chobot},
Stefania\indeksosoby[Poziemska]{Stefania}{Osica}{1897}{Chobot} oraz
Michał\indeksosoby{Michał}{Poziemski}{1901}{Chobot}.
Franciszek\indeksosoby{Franciszek}{Poziemski}{1894}{Chobot} zmarł
w~1917 roku jako kawaler w~wieku 23~lat, natomiast jego najmłodsza siostra,
Czesława\indeksosoby{Czesława}{Poziemska}{1906}{Desno}, odeszła w~1922~roku,
dożywając zaledwie szesnastu lat.

\begin{figure}[!ht]
    \centering \includegraphics[width=1.0\linewidth]{
        1917_Franciszek_Poziemski_akt_zgonu_parafia_Długa_Kościelna_wpis_51.jpg}
    \captionsetup{format=hang}
    \caption{Akt zgonu Franciszka Poziemskiego -- par. Długa Kościelna
    1917~rok (51/1917) \cite{par_dluga}.}
    \label{fig:fpoziemski_1917}
\end{figure}

Dnia 9~lutego 1919~roku Aleksander
Poziemski\indeksosoby{Aleksander}{Poziemski}{1892}{Chobot} wziął ślub
w~parafii świętej Anny w~Długiej Kościelnej z~Franciszką
Baliszewską\indeksosoby[Baliszewska]{Franciszka}{Poziemska}{1895}{Chobot}
(ur.~1895~r.~-- zm.~1968~r.), dwudziestotrzyletnią panną urodzoną
i~zamieszkałą w~Chobocie, córką Onufrego i~Petroneli (z~domu Jackiewicz)
Baliszewskich (patrz ryc. \ref{fig:apoziemski_1919}). Para po~ślubie
zamieszkała w~rodzinnej wsi Aleksandra, gdzie dochowała się wspólnie
co~najmniej dwójki dzieci:
Kazimiery\indeksosoby[Poziemska][Jurkowska]{Kazimiera}{Andrychowska}{1920}{Desno}
(ur.~1920~r.~-- zm.~1997~r.) oraz
Stanisława\indeksosoby{Stanisław}{Poziemski}{1928}{Desno} (ur.~1928~r.~--
zm.~1930~r.). Co~ciekawe małżonkowie Poziemscy w~pewnym momencie swojego życia
wyprowadzili się z~~Desna w~okolice miasta Morąg leżącego w~województwie
warmińsko-mazurskim, być może ze~względu na~fakt, iż~ich córka Kazimiera
w~1962~roku wzięła tam ślub z~Wiktorem Andryszewskim\footnote{Był to~jej~drugi
ślub -- jej poprzedni mąż Marian Jurkowski zginął w~1945~roku w~obozie
koncentracyjnym Mauthausen.}.
Aleksander\indeksosoby{Aleksander}{Poziemski}{1892}{Chobot}
i~Franciszka\indeksosoby[Baliszewska]{Franciszka}{Poziemska}{1895}{Chobot}
Poziemscy dożyli tam już swoich dni -- zostali pochowani na~cmentarzu
komunalnym w~Antoniewie, w~gminie Morąg (patrz ryc. \ref{fig:apoziemski_1968}).

\begin{figure}[!ht]
    \centering \includegraphics[width=0.93\linewidth]{
        1919_Aleksander_Poziemski_Franciszka_Baliszewska_akt_ślubu_parafia_Długa_Kościelna_wpis_6.jpg}
    \captionsetup{format=hang}
    \caption{Akt ślubu Aleksandra Poziemskiego i~Franciszki Baliszewskiej --
    par. Długa Kościelna 1919~rok (6/1919) \cite{par_dluga}.}
    \label{fig:apoziemski_1919}
\end{figure}

\begin{figure}[!ht]
    \centering \includegraphics[width=0.9\linewidth]{
        1968_Aleksander_Poziemski_zdjęcie_nagrobka1_cmentarz_Antoniewo.jpg}
    \captionsetup{format=hang}
    \caption{Zdjęcie grobu małżonków Aleksandra oraz Franciszki Poziemskich,
    znajdującego się na cmentarzu komunalnym w~Antoniewie, gmina Morąg
    (sektor: 5, numer: 165)}
    \label{fig:apoziemski_1968}
\end{figure}

Dnia 19~lipca 1931~roku trzydziestoczteroletnia Stefania
Poziemska\indeksosoby[Poziemska]{Stefania}{Osica}{1897}{Chobot} wyszła
za~mąż w~parafii świętej Anny w~Długiej Kościelnej za~Jana
Osicę\indeksosoby{Jan}{Osica}{1903}{Dłużew} (ur.~1903~r~-- zm.~?)
dwudziestoośmioletniego kawalera pochodzącego ze~wsi Dłużew (w~gminie Siennica),
syna Aleksandra i~Cecyli (z~domu Pazio) Osiców. Para prawdopodobnie nie
dochowała się potomstwa.

\begin{figure}[!ht]
    \centering \includegraphics[width=0.86\linewidth]{
        1903_Jan_Osica_akt_chrztu_parafia_Kołbiel_wpis_184.jpg}
    \captionsetup{format=hang}
    \caption{Akt chrztu Jana Osicy -- par. Kołbiel 1903~rok (184/1903)
    \cite{par_kolbiel2}.}
    \label{fig:josica_1903}
\end{figure}

\begin{figure}[!ht]
    \centering \includegraphics[width=0.86\linewidth]{
        1895_Magdalena_Gomulska_Zawadzka_Poziemska_akt_chrztu_parafia_Długa_Kościelna_wpis_66.jpg}
    \captionsetup{format=hang}
    \caption{Akt chrztu Magdaleny Poziemskiej z~domu Zawadzkiej -- par.
    Długa Kościelna 1895~rok (66/1895) \cite{par_dluga}.}
    \label{fig:mzawadzka_1896}
\end{figure}

Dnia 5~maja 1929~roku dwudziestosiedmioletni Michał
Poziemski\indeksosoby{Michał}{Poziemski}{1901}{Chobot} wziął ślub
ze~straszą od~niego o~sześć lat Magdaleną
Gomulską\indeksosoby[Zawadzka][Gomulska]{Magdalena}{Poziemska}{1895}{Cisie}
z~domu Zawadzką (ur.~1895~r~-- zm.~1970~r.), wdową po~zmarłym Antonim
Gomulskim\indeksosoby{Antoni}{Gomulski}{1882}{Mikanów}, córką Marcina
i~Marianny (z~domu Zawadzkiej) Zawadzkich. W~momemcie ich ślubu Magdalena
wychowywała już dwóch synów: Marjana
Zawadzkiego\indeksosoby{Marjan}{Zawadzki}{1917}{Desno} (ur.~1917~r.~-- zm.~?)
-- nieślubnego syna Stefana
Gomulskiego\indeksosoby{Stefan}{Gomulski}{1896}{Desno} oraz Romana
Gomulskiego\indeksosoby{Roman}{Gomulski}{1922}{Desno} (ur.~1922~r.~--
zm.~1976~r.) -- syna swojego pierwszego męża\footnote{O~czym szerzej pisałem
w~rozdziale V~niniejszej książki}. Michał\indeksosoby{Michał}{Poziemski}{1901}{Chobot}
i~Magdalena\indeksosoby[Zawadzka][Gomulska]{Magdalena}{Poziemska}{1895}{Cisie}
Poziemscy przez 18~lat małżeństwa doczekali się wspólnie czwórki
dzieci\footnote{Z~czego tylko ich najmłodsza córka
Jadwiga\indeksosoby[Poziemska]{Jadwiga Leokadia}{Gomulska}{1935}{Desno}
przeżyła dłużej niż kilka miesięcy -- dożyła wieku 79~lat dochowując się
wspólnie z~mężem Stanisławem
Gomulskim\indeksosoby{Stanisław}{Gomulski}{1926}{Desno} (ur.~1926~r.~--
zm.~1923~r.) czwórki dzieci, siedmiorga wnucząt i~jedenaściorga prawnucząt.}:

\begin{itemize}
    \item Marianny\indeksosoby{Marianna}{Poziemska}{1931}{Desno}
    (ur.~1931~r.~-- zm.~1931~r.),
    \item Honoraty\indeksosoby{Honorata}{Poziemska}{1931}{Desno}
    (ur.~1931~r.~-- zm.~1931~r.),
    \item Józefa\indeksosoby{Józefa}{Poziemska}{1933}{Desno}
    (ur.~1933~r.~-- zm.~1933~r.),
    \item Jadwigii Leokadii\indeksosoby[Poziemska]{Jadwiga
    Leokadia}{Gomulska}{1935}{Desno}\footnote{Jadwiga Leokadia Poziemska była
    babcią autora niniejszej książki.} (ur.~1935~r.~-- zm.~2014~r.).
\end{itemize}

\begin{figure}[!ht]
    \vspace*{0.4cm}
    \centering \includegraphics[width=1.0\linewidth]{
        1947_Michał_Poziemski_zdjęcie_nagrobka1_cmentarz_Długa_Kościelna.jpg}
    \captionsetup{format=hang}
    \caption{Zdjęcie grobu małżonków Michała oraz Magdaleny Poziemskich,
    znajdującego się na cmentarzu parafialnym parafii św. Anny w~Długiej
    Kościelnej (sektor: A5, rząd: A, numer: 16)}
    \label{fig:mpoziemski_1947}
\end{figure}

\newpage
\begin{figure}[!ht]
\centering
\begin{tikzpicture}[
  % Definiujemy stałe odstępy
  node distance=1cm and 1cm, 
  every node/.style={draw, rectangle, align=center, minimum width=3.8cm,
  inner sep=5pt}]

% --- POKOLENIE 0 ---
\node (Osiowa0) [draw=none, minimum width=0pt, inner sep=0pt] at (0,0) {};
\node (Franciszek) [left=1.5cm of Osiowa0] {Franciszek Poziemski \\ {
    \footnotesize \emph{?, ? r.}}};
\node (Anna) [right=1.5cm of Osiowa0] {Anna Borakowska \\ {
    \footnotesize \emph{?, ? r.}}};

\draw[-] (Franciszek) -- (Anna)
    node[midway, below, draw=none, minimum width=0pt,
    font=\scriptsize] (Slub0) {\emph{Kamionna, 1801 r.}};

% --- POKOLENIE 1 ---
\node (Osiowa1) [draw=none, minimum width=0pt, inner sep=0pt,
below=1.5cm of Slub0] {};
\node (Andrzej) [left=1.5cm of Osiowa1] {Andrzej Poziemski \\ {
    \footnotesize \emph{Budy Łosiewickie, 1806 r.}}};
\node (Rozalia) [right=1.5cm of Osiowa1] {Rozalia Rosik \\ {
    \footnotesize \emph{Ossów, 1800 r.}}};

\draw[-] (Andrzej) -- (Rozalia)
    node[midway, above, draw=none, minimum width=0pt,
    font=\scriptsize] (Slub1_Gora) { }
    node[midway, below, draw=none, minimum width=0pt,
    font=\scriptsize] (Slub1) {\emph{Okuniew, 1828 r.}};

\draw[->] (Slub0) -- (Slub1_Gora);

% --- POKOLENIE 2 ---
\node (Osiowa2) [draw=none, minimum width=0pt, inner sep=0pt,
below=1.5cm of Slub1] {};
\node (Franciszek2) [left=1.5cm of Osiowa2] {Franciszek Poziemski \\ {
    \footnotesize \emph{Budziska, 1829 r.}}};
\node (Magdalena) [right=1.5cm of Osiowa2] {Magdalena Grzybowska \\ {
    \footnotesize \emph{Kaleń, 1827 r.}}};
    
\draw[-] (Franciszek2) -- (Magdalena)
    node[midway, above, draw=none, minimum width=0pt,
    font=\scriptsize] (Slub2_Gora) { }
    node[midway, below, draw=none, minimum width=0pt,
    font=\scriptsize] (Slub2) {\emph{Okuniew, 1851 r.}};

\draw[->] (Slub1) -- (Slub2_Gora);

% --- POKOLENIE 3 ---
\node (Osiowa3) [draw=none, minimum width=0pt, inner sep=0pt,
below=1.5cm of  Slub2] {};
\node (Antoni) [left=1.5cm of Osiowa3] {Antoni Poziemski \\ {
    \footnotesize \emph{Cisie, 1864 r.}}};
\node (Józefa) [right=1.5cm of Osiowa3] {Józefa Tomaszewska \\ {
    \footnotesize \emph{Chobot, 1864 r.}}};

\draw[-] (Antoni) -- (Józefa)
    node[midway, above, draw=none, minimum width=0pt,
    font=\scriptsize] (Slub3_Gora) { }
    node[midway, below, draw=none, minimum width=0pt,
    font=\scriptsize] (Slub3) {\emph{Długa Kościelna, 1886 r.}};

\draw[->] (Slub2) -- (Slub3_Gora);

% --- POKOLENIE 4 ---
\node (Osiowa4) [draw=none, minimum width=0pt, inner sep=0pt,
below=1.5cm of  Slub3] {};
\node (Michał) [left=1.5cm of Osiowa4] {Michał Poziemski \\
{\footnotesize \emph{Chobot, 1901 r.}}};
\node (Magdalena2) [right=1.5cm of Osiowa4] {Magdalena Gomulska \\
{\footnotesize \emph{Cisie, 1895 r.}}};

\draw[-] (Michał) -- (Magdalena2)
    node[midway, above, draw=none, minimum width=0pt,
    font=\scriptsize] (Slub4_Gora) { }
    node[midway, below, draw=none, minimum width=0pt,
    font=\scriptsize] (Slub4) {\emph{Długa Kościelna, 1929 r.}};

\draw[->] (Slub3) -- (Slub4_Gora);

% --- POKOLENIE 5 ---
\node (Jadwiga) [below=1.5cm of Slub4] {Jadwiga Leokadia Poziemska \\
{\footnotesize \emph{Desno, 1935 r.}}};

\draw[->] (Slub4) -- (Jadwiga);

\end{tikzpicture}
\end{figure}